\chapter{Domain and Mesh Representation}
In this subsection, we present definitions of non-overlapping and overlapping domain decompositions. For overlapping decompositions, we introduce several variants, including minimal overlapping, uniform overlapping, and nonuniform overlapping, which differ in how the subdomains extend beyond their base regions. These definitions are provided for both geometric (non-mesh-aligned) and mesh-conforming decompositions. We also highlight alternative definitions and approaches that may appear in the literature.

\section{Geometric Decomposition}

This subsection considers the geometric decomposition of a bounded domain $\Omega \subset \mathbb{R}^{d}$, where $\Omega$ is a non-empty, connected, open, and bounded subset of $\mathbb{R}^{d}$. While $\Omega$ may have a smooth or curved boundary $\partial\Omega$, we assume that $\partial\Omega$ satisfies the strong local Lipschitz condition so that standard Sobolev space and trace results remain valid (see, e.g., Adams and Fournier~\cite{AdamsFournier2003}). We first introduce a non-overlapping decomposition of such a domain and then proceed to the overlapping decomposition.

\begin{definition}
Let $\Omega \subset \mathbb{R}^d$ be a bounded domain. A \textit{non-overlapping decomposition} of $\Omega$ is a finite collection of subdomains $\{\Omega_{i}\}$ such that
\begin{equation*}
    \Omega_{i} \cap \Omega_{j} = \emptyset \quad\text{ for } i \neq j,
\end{equation*}
and
\begin{equation*}
    \overline{\Omega} = \bigcup\overline{\Omega}_{i}.
\end{equation*}
For each subdomain $\Omega_i$, let
\begin{equation*}
    \Gamma_{i} := \partial \Omega_i
\end{equation*}
denote its boundary. The boundary can be decomposed into two disjoint parts, 
\begin{equation*}
    \Gamma_{i} = \Gamma_{i}^{\text{ext}} \cup \Gamma_{i}^{\text{int}},
\end{equation*}
where
\begin{equation*}
    \Gamma_{i}^{\text{ext}} := \partial \Omega_i \cap \partial \Omega
\end{equation*}
is the \textit{exterior boundary}, and
\begin{equation*}
    \Gamma_{i}^{\text{int}} := \partial \Omega_{i} \setminus \partial\Omega
\end{equation*}
is the \textit{interface boundary}, which is shared with neighboring subdomains.
\end{definition}
\noindent After introducing a non-overlapping decomposition, it is useful to classify subdomains according to whether they touch the boundary of the global domain. This distinction is important in domain decomposition methods, as boundary subdomains are adjacent to the global boundary $\partial \Omega$, whereas interior subdomains are entirely surrounded by other subdomains. Formally, we define boundary and interior subdomains as follows.
\begin{definition}
Let $\{\Omega_i\}_{i=1}^N$ be a non-overlapping decomposition of a bounded domain $\Omega \subset \mathbb{R}^d$. A subdomain $\Omega_i$ is called a \textit{boundary subdomain} if its exterior boundary is non-empty,
\begin{equation*}
    \Gamma_i^{\text{ext}} \neq \emptyset,
\end{equation*}
and it is called an \textit{interior subdomain} if its exterior boundary is empty,
\begin{equation*}
    \Gamma_i^{\text{ext}} = \emptyset.
\end{equation*}
\end{definition}
\noindent Each subdomain $\Omega_i$ has a boundary $\Gamma_i = \partial \Omega_i$, which can be split into the exterior boundary $\Gamma_i^{\text{ext}}$ and the interface (interior) boundary $\Gamma_i^{\text{int}}$.  
The exterior boundary indicates whether the subdomain touches the boundary of the global domain $\Omega$, while the interface boundary corresponds to the portion of $\partial \Omega_i$ shared with neighboring subdomains. The union of all interface boundaries across the decomposition forms the \textbf{interface} (or \textbf{skeleton}) of the domain decomposition (see Brenner and Scott~\cite{BrennerScott2008}). This interface is central in domain decomposition methods, as it represents the region where subdomains interact and exchange information.
\begin{definition}
Let $\{\Omega_i\}_{i=1}^N$ be a non-overlapping decomposition of $\Omega$. The \textit{interface} (also called the \textit{skeleton}) of the decomposition is defined as
\begin{equation*}
    \Gamma^{\text{int}} := \bigcup\Gamma_{i}^{\text{int}}.
\end{equation*}
\end{definition}
\noindent While non-overlapping decompositions partition the domain into disjoint subdomains, it is often advantageous in numerical methods to consider subdomains that may partially overlap. Overlapping subdomains allow regions to be shared between neighboring subdomains, facilitating the transfer of information across interfaces and improving the convergence and stability of iterative algorithms. To make this concept precise, we provide the following formal definition of an overlapping decomposition.

\begin{definition}
Let $\Omega \subset \mathbb{R}^d$ be a bounded domain. An \textit{overlapping decomposition} of $\Omega$ is a finite collection of subdomains $\{\Omega_{i}\}$ such that
\begin{equation*}
    \Omega = \bigcup\Omega_{i}.
\end{equation*}
For each subdomain $\Omega_i$, let
\begin{equation*}
    \Gamma_{i} := \partial \Omega_i
\end{equation*}
denote its boundary. The boundary can be decomposed into two disjoint parts, 
\begin{equation*}
    \Gamma_{i} = \Gamma_{i}^{\text{ext}} \cup \Gamma_{i}^{\text{int}},
\end{equation*}
where
\begin{equation*}
    \Gamma_{i}^{\text{ext}} := \partial \Omega_i \cap \partial \Omega
\end{equation*}
is the \textbf{exterior boundary}, and
\begin{equation*}
    \Gamma_{i}^{\text{int}} := \partial \Omega_{i} \setminus \partial\Omega
\end{equation*}
is the \textbf{interface boundary}, which is shared with neighboring subdomains.
\end{definition}
\noindent Let us now define what we mean by "neighboring subdomains" and provide an explicit definition.
\begin{definition}
Let $\{\Omega_i\}_{i=1}^N$ be a decomposition of a bounded domain $\Omega \subset \mathbb{R}^d$. For $i \neq j$, define 
\begin{equation*}
    \Gamma_{ij} := \Gamma_i^{\text{int}} \cap \overline{\Omega}_j.
\end{equation*}
We say that $\Omega_i$ and $\Omega_j$ are \textbf{neighboring subdomains} if
\begin{equation*}
    \Gamma_{ij} \neq \emptyset.
\end{equation*}
We denote by $\mathcal{N}_i$ the set of indices of all subdomains neighboring $\Omega_{i}$, i.e.,
\begin{equation*}
    \mathcal{N}_i := \{ j \in \{1, \dots, N\} \setminus \{i\} : \Gamma_{ij} \neq \emptyset \}.
\end{equation*}
\end{definition}
This definition holds for both non-overlapping and overlapping decompositions, but it is more relevant in the context of overlapping decompositions, 
where subdomains can share regions. In non-overlapping decompositions, neighboring subdomains typically only touch at their boundaries, while in 
overlapping decompositions, they can have a non-empty intersection. Also, notice that the interface boundary of a subdomain can be expressed as
\begin{equation*}
    \Gamma_{i}^{\text{int}} = \bigcup_{j \in \mathcal{N}_i} \Gamma_{ij},
\end{equation*}
where the union is \textbf{not necessarily disjoint}. 

\vspace{8pt}
\noindent The standard definition of an overlapping decomposition of a domain is quite general and lacks practical clarity. We therefore adopt a definition based on non-overlapping decomposition, which is typically used in applications. In many domain decomposition methods, it is convenient to extend non-overlapping subdomains by a small amount to create overlapping subdomains.  
This extension allows neighboring subdomains to communicate information across their interfaces, which is important for the convergence of iterative methods.  
In the most general setting, the amount of extension may vary across subdomains and even depend on position within the subdomain. The following definition introduces this general case.

\begin{definition}{Overlapping Domain Decomposition}{overlapping_decomposition}
Let $\{\widetilde{\Omega}_i\}_{i=1}^N$ be a non-overlapping decomposition of a bounded domain $\Omega\subset \mathbb{R}^{d}$. An \textbf{overlapping decomposition} of $\Omega$ is a finite collection of subdomains $\{\Omega_i\}_{i=1}^N$ defined by
\begin{equation*}
    \Omega_i := \{ x \in \Omega : \operatorname{dist}(x, \widetilde{\Omega}_i) < \delta_i(x) \},
\end{equation*}
where $\delta_i: \Omega \to \mathbb{R}^{+}$ is a bounded, continuous function, called the \textit{extension function} for the $i$-th subdomain.
\end{definition}

In the definition of neighboring subdomains, we define $\Gamma_{ij}$, which is the interface boundary of a subdomain $\Omega_i$ that is shared with a neighboring 
subdomain $\Omega_j$. In the case of overlapping decomposition, $\Gamma_{ij}$ may be defined differently, as it can be seen in the literature. There are two common 
approaches to defining the interface boundary between subdomains in overlapping decompositions: Let $\{\widetilde{\Omega}_i\}_{i=1}^N$ be a non-overlapping decomposition of a bounded 
domain $\Omega \subset \mathbb{R}^d$ and $\{\Omega_i\}_{i=1}^N$ be an overlapping decomposition of $\Omega$ obtained by extending the non-overlapping subdomains 
$\widetilde{\Omega}_i$, then for $i \neq j$, the interface boundary $\Gamma_{ij}$ can be defined in two ways:
\begin{enumerate}[label=(\arabic*)]
    \item $\Gamma_{ij} := \Gamma_{i}^{\text{int}} \cap \overline{\Omega}_j$
    \item $\Gamma_{ij} := \Gamma_{i}^{\text{int}} \cap \overline{\widetilde{\Omega}}_j$
\end{enumerate}
Making note of this distinction is important, as it affects the analysis and implementation of domain decomposition methods. The first definition considers the 
interface boundary as the portion of $\Gamma_i^{\text{int}}$ that intersects with the extended subdomain $\Omega_j$, while the second definition restricts the 
interface boundary to the portion that intersects with the original non-overlapping subdomain $\widetilde{\Omega}_j$. The choice between these definitions can 
influence the properties of the interface and the communication between subdomains in iterative methods, as well as the theoretical analysis of convergence and 
stability. Let us call $\Gamma_{ij}^{(1)}$ and $\Gamma_{ij}^{(2)}$ the interface boundaries defined by the first and second approaches, respectively. The interface 
boundary $\Gamma_i^{\text{int}}$ can be expressed as the union of $\Gamma_{ij}$ over all neighboring subdomains, independent of $\Gamma_{ij}$ definition, that is,
\begin{equation*}
    \Gamma_i^{\text{int}} = \bigcup_{j \in \mathcal{N}_i} \Gamma_{ij}^{(1)} = \bigcup_{j \in \mathcal{N}_i} \Gamma_{ij}^{(2)},
\end{equation*}  
and neighboring subdomains are exactly the same under both definitions, i.e., $\mathcal{N}_i^{(1)} = \mathcal{N}_i^{(2)}$. Let us now analyze the intersection 
of interface boundaries between neighboring subdomains under both definitions. For $j \neq k$ and $j, k \in \mathcal{N}_i$, we have
\begin{enumerate}[label = (\arabic*)]
    \item If $j\in \mathcal{N}_{k}$ or $k \in \mathcal{N}_{j}$ (subdomains $\Omega_j$ and $\Omega_k$ are neighbors), then
    \begin{equation*}
        \dim\left(\Gamma_{ij}^{(1)} \cap \Gamma_{ik}^{(1)}\right) = \dim\left(\Gamma_{i}^{\text{int}} \cap \overline{\Omega}_j \cap \overline{\Omega}_k\right) = d-1.
    \end{equation*}
    Otherwise, if $j \notin \mathcal{N}_{k}$ and $k \notin \mathcal{N}_{j}$ (subdomains $\Omega_j$ and $\Omega_k$ are not neighbors), then $\Gamma_{ij}^{(1)} \cap \Gamma_{ik}^{(1)} = \emptyset$, and thus\footnote{In the topological sense, the empty set has dimension $-1$.}
    \begin{equation*}
        \dim\left(\Gamma_{ij}^{(1)} \cap \Gamma_{ik}^{(1)}\right) = -1.
    \end{equation*}
    \item If $j\in \mathcal{N}_{k}$ or $k \in \mathcal{N}_{j}$ (subdomains $\Omega_j$ and $\Omega_k$ are neighbors), then
    \begin{equation*}
        \dim\left(\Gamma_{ij}^{(2)} \cap \Gamma_{ik}^{(2)}\right) = \dim\left(\Gamma_{i}^{\text{int}} \cap \overline{\widetilde{\Omega}}_j \cap \overline{\widetilde{\Omega}}_k\right) = d-2.
    \end{equation*}
    Otherwise, if $j \notin \mathcal{N}_{k}$ and $k \notin \mathcal{N}_{j}$ (subdomains $\Omega_j$ and $\Omega_k$ are not neighbors), then $\Gamma_{ij}^{(2)} \cap \Gamma_{ik}^{(2)} = \emptyset$, and thus
    \begin{equation*}
        \dim\left(\Gamma_{ij}^{(2)} \cap \Gamma_{ik}^{(2)}\right) = -1.
    \end{equation*}
\end{enumerate}
This observation is crucial as it shows that the choice of interface boundary definition affects the geometric properties of the interfaces between subdomains. 
The first definition allows for higher-dimensional intersections between interface boundaries of neighboring subdomains, while the second definition restricts 
these intersections to lower-dimensional sets. This has implications for the analysis of domain decomposition methods, particularly in terms of convergence rates 
and stability, as well as for the implementation of communication between subdomains in parallel computing environments. In the following subsection, we will see
the importance of the choice of the interface boundary decomposition.

\noindent A natural special case of the general overlapping decomposition arises when the extension function is constant within each subdomain, 
but may still vary from one subdomain to another. In this case, each subdomain is enlarged by a fixed distance $\delta_i$, leading to what we call a \textit{nonuniform overlapping decomposition}. The following definition formalizes this setting.

\begin{definition}
Let $\{\widetilde{\Omega}_i\}_{i=1}^N$ be a non-overlapping decomposition of a bounded domain $\Omega \subset \mathbb{R}^d$. A \textbf{nonuniform overlapping decomposition} of $\Omega$ is a finite collection of subdomains $\{\Omega_{i}^{\delta_{i}}\}_{i=1}^N$ defined by
\begin{equation*}
    \Omega_i^{\delta_{i}} := \{ x \in \Omega : \operatorname{dist}(x, \widetilde{\Omega}_i) < \delta_i \},
\end{equation*}
where $\delta_i > 0$ is the \textit{extension distance} for the $i$-th subdomain.
\end{definition}
\noindent The distance $\operatorname{dist}(x,\widetilde{\Omega}_{i})$ is defined by
\begin{equation*}
    \operatorname{dist}(x,\widetilde{\Omega}_{i}) := \inf_{y\in\widetilde{\Omega}_{i}} \|x-y\|,
\end{equation*}
where $\|\cdot\|$ denotes a norm on $\mathbb{R}^{d}$. The choice of norm affects the geometry of the extended subdomains. For instance, using the Euclidean norm $\ell^2$ enlarges each subdomain with a ball, producing rounded overlaps, whereas using the maximum norm $\ell^\infty$ enlarges subdomains with axis-aligned boxes, resulting in box-shaped overlaps. This is illustrated in Figure~\ref{fig1}.
\begin{figure}[H]
\centering
\begin{subfigure}[b]{0.2\textwidth}
    \centering
    \includegraphics[width=\textwidth]{img1.jpg}
    \caption{Euclidean norm $\ell^2$}
    \label{fig11}
\end{subfigure}
\hspace{20pt}
\begin{subfigure}[b]{0.2\textwidth}
    \centering
    \includegraphics[width=\textwidth]{img2.jpg}
    \caption{Maximum norm $\ell^\infty$}
    \label{fig12}
\end{subfigure}
\caption{Comparison of uniform overlapping subdomains under different norms.}
\label{fig1}
\end{figure}

\noindent By definition, every point in $\Omega_i^{\delta_i}$ lies within distance $\delta_i$ of $\widetilde{\Omega}_i$, that is, 
\begin{equation*}
    \sup_{x \in \Omega_i^{\delta_{i}}} \operatorname{dist}(x, \widetilde{\Omega}_i) \leq \delta_i.
\end{equation*}
Moreover, since $\widetilde{\Omega}_i \subset \Omega_i^{\delta_{i}}$, we have 
\begin{equation*}
    \sup_{y \in \widetilde{\Omega}_i} \operatorname{dist}(y, \Omega_i^{\delta_{i}}) = 0.
\end{equation*}
It follows that the Hausdorff distance between $\Omega_i^{\delta_i}$ and $\widetilde{\Omega}_i$ satisfies
\begin{equation*}
    d_H(\Omega_i^{\delta_{i}}, \widetilde{\Omega}_i) = \max\left\{ \sup_{x \in \Omega_i^{\delta_{i}}} \operatorname{dist}(x, \widetilde{\Omega}_i), \sup_{y \in \widetilde{\Omega}_i} \operatorname{dist}(y, \Omega_i^{\delta_{i}}) \right\} \leq \delta_i,
\end{equation*}
where we consider the metric space $(M,d)$ with $M = \Omega$ and $d$ induced by the norm used in $\operatorname{dist}$. In other words, the extended subdomain $\Omega_i^{\delta_{i}}$ lies entirely within a $\delta_i$-neighborhood of $\widetilde{\Omega}_i$, while $\widetilde{\Omega}_i$ is fully contained in $\Omega_i^{\delta_{i}}$. This shows rigorously that the Hausdorff distance between the subdomain and its $\delta_i$-extension is at most $\delta_i$.

\noindent In general, the extension distance may vary across subdomains, denoted by $\delta_i$ for the $i$-th subdomain. In much of the literature, however, the extension distance is taken to be the same for all subdomains, i.e., $\delta_i = \delta$ for all $i$. The resulting decomposition is then called a \textit{uniform overlapping decomposition}. In this case, each overlapping subdomain $\Omega_i^{\delta}$ is obtained by enlarging its corresponding non-overlapping subdomain $\widetilde{\Omega}_i$ by the fixed distance $\delta$, so that
\begin{equation*}
    d_H(\widetilde{\Omega}_i, \Omega_i^{\delta}) \leq \delta.
\end{equation*}
Although $\delta$ is sometimes referred to as the \textit{overlap width}, it geometrically represents the outward extension of each subdomain. Since the overlap between neighboring subdomains results from two such extensions, $\delta$ corresponds to half of the total overlap. For this reason, we refer to $\delta$ as the \textit{uniform extension distance}.
\begin{definition}{Uniform Overlapping Domain Decomposition}{uniform_overlapping_decomposition}
Let $\{\widetilde{\Omega}_i\}_{i=1}^N$ be a non-overlapping decomposition of a bounded domain $\Omega \subset \mathbb{R}^d$. A \textbf{uniform overlapping decomposition} of $\Omega$ is a finite collection of subdomains $\{\Omega_i^{\delta}\}_{i=1}^N$ defined by
\begin{equation*}
    \Omega_i^{\delta} := \{ x \in \Omega : \operatorname{dist}(x, \widetilde{\Omega}_i) < \delta \},
\end{equation*}
where $\delta > 0$ is the \textit{uniform extension distance}.
\end{definition}
\noindent In practical applications, it is common to impose a minimal overlap between subdomains to ensure sufficient communication across interfaces. In the work of Gander and Zhao~\cite{GanderZhao2002}, the domain $\Omega$ is decomposed into $N$ nonuniform overlapping subdomains $\{\Omega_i\}_{i=1}^N$, where each subdomain is defined by
\begin{equation*}
    \Omega_i = \{ x \in \Omega : \operatorname{dist}(x, \widetilde{\Omega}_i) < \delta_i(x) \}, \qquad \delta_i(x) \geq \delta > 0, \quad x \in \Omega \setminus \widetilde{\Omega}_i,
\end{equation*}
with $\widetilde{\Gamma}_i = \partial \widetilde{\Omega}_i$ and $\Gamma_i = \partial \Omega_i$ denoting the boundaries of the non-overlapping and overlapping subdomains, respectively. They further assume that the distance between the interface boundaries satisfies
\begin{equation*}
    \operatorname{dist}(\widetilde{\Gamma}_i^{\text{int}}, \Gamma_i^{\text{int}}) \geq \delta,
\end{equation*}
so that each overlapping subdomain extends at least a distance $\delta$ beyond its corresponding non-overlapping subdomain (see Figure~\ref{fig2}). This assumption is important because, in general, a $\delta$-extension as defined above does not automatically guarantee that the Hausdorff distance between the original and extended subdomain (or their boundaries) is at least $\delta$ everywhere. Geometric constraints, such as the shape of the global domain or the arrangement of neighboring subdomains, may prevent certain regions from achieving the full extension. By explicitly assuming a minimal extension, one ensures that all subdomains overlap sufficiently, which is essential for the convergence and stability of domain decomposition methods.
\begin{figure}[H]
    \centering
    \includegraphics[width=0.3\textwidth]{delta.jpg}
    \caption{Decomposition into overlapping subdomains such that an extension of $\delta$ is guaranteed.~\cite{GanderZhao2002}}
    \label{fig2}
\end{figure}
\begin{definition}
Let $\{\widetilde{\Omega}_i\}_{i=1}^N$ be a non-overlapping decomposition of a bounded domain $\Omega \subset \mathbb{R}^d$. A collection of overlapping subdomains 
$\{\Omega_i\}_{i=1}^N$ defined by
\begin{equation*}
    \Omega_i := \{ x \in \Omega : \operatorname{dist}(x, \widetilde{\Omega}_i) < \delta_i(x) \},
\end{equation*}
where the extension function $\delta_i : \Omega \to (0,\infty)$ is a bounded, continuous function, is called a \textbf{overlapping decomposition} 
of $\Omega$ with $\delta$-extension if the \textit{minimal extension condition}
\begin{equation*}
    \operatorname{dist}(\widetilde{\Gamma}_i^{\text{int}}, \Gamma_i^{\text{int}}) \geq \delta,
\end{equation*}
is satisfied for some constant $\delta > 0$, called the \textit{minimal extension distance}.
\end{definition}
\begin{remark}
The term "minimal" does not mean that the overlap is the smallest possible. It indicates that each subdomain is extended by at least the prescribed minimal 
extension distance $\delta$ beyond its non-overlapping boundary. As a result, the minimal overlap distance, defined as the smallest overlap among all neighboring 
subdomain pairs, is at least $2\delta$. 

\vspace{6pt}
\noindent For each extended subdomain $\Omega_i$ with minimal extension distance $\delta$, the extension distances satisfy
\begin{equation*}
    \delta_i(x) \ge \operatorname{dist}(\widetilde{\Gamma}_i^{\text{int}}, \Gamma_i^{\text{int}}) \geq \delta>0, 
    \qquad 1 \leq i \leq N, \quad x \in \Omega \setminus \widetilde{\Omega}_i.
\end{equation*}
Geometrically, this means that each subdomain $\Omega_i$ contains a uniform neighborhood of width at least $\delta$ around the interior boundary of the 
corresponding subdomain $\widetilde{\Omega}_i$, guaranteeing a minimal overlap between neighboring subdomains (see Figure~\ref{fig2}).
\end{remark}

\section{Mesh-Conforming Decomposition}

Typically, theoretical results assume that $\Omega$ is a convex polytopal domain\footnote{There are two standard definitions of a polytope, and in both cases, the notion of a convex polytope is equivalent.} (convex polygonal in $\mathbb{R}^{2}$, convex polyhedral in $\mathbb{R}^{3}$). Convexity is not essential, but it simplifies analysis and guarantees that certain inequalities and estimates hold with sharp constants. For example, the Poincaré constant can be made explicit for convex domains (see Bebendorf~\cite{Bebendorf2003Poincare}). One might think that restricting to polytopal domains is limiting. In practice, however, domains with smooth or curved boundaries are often approximated by polygonal (2D) or polyhedral (3D) domains. The resulting error, called the isoparametric error, can be rigorously quantified and is considered one of the variational crimes. To avoid these issues, we assume in this paper that $\Omega \subset \mathbb{R}^{d}$ is a polytopal domain, irrespective of convexity. Note that under this assumption, the strong local Lipschitz condition is automatically satisfied, ensuring that standard Sobolev space and trace results hold.

Assuming that $\Omega \subset \mathbb{R}^{d}$ is a polytopal domain, the next step is to construct a non-overlapping decomposition of the domain, i.e., to find a finite collection of polytopal subdomains $\{\widetilde{\Omega}_{i}\}$ satisfying the conditions given in the definition above. Since we are employing the finite element method (FEM), we first generate a finite element subdivision (mesh) of the domain $\Omega$, consisting of element domains. These element domains are then grouped to form subdomains satisfying the conditions of the non-overlapping domain decomposition. Once the non-overlapping domain decomposition is constructed, an overlapping decomposition can be obtained by adding all neighbouring elements\footnote{Here, we assume vertex neighbourhoods rather than edge neighbourhoods.} to each subdomain $\widetilde{\Omega}_{i}$. Since the domain decomposition is not unique, care must be taken: one may include more elements in certain directions or extend uniformly, which affects the extension distance and, consequently, the theoretical bounds derived in the subsequent analysis.

To define a mesh-conforming decomposition of a polytopal domain, we first consider the cases $d = 1, 2$ and then extend the definition to the general case.

\begin{enumerate}[label=(\arabic*)]
\item If $d = 1$, then the domain $\Omega = (a,b)$ is an open interval. A non-overlapping decomposition of $\Omega$ into $N$ subdomains 
consists of subintervals 
\begin{equation*}
    \widetilde{\Omega}_i = (x_{i-1}, x_i),
\end{equation*}
where $x_0 = a$, $x_N = b$, and $x_i > x_{i-1}$ for $i = 1, \ldots, N$.
The partition may be uniform or non-uniform. In the uniform case,
\begin{equation*}
    x_i - x_{i-1} = h \quad \text{for } i = 1, \ldots, N,
\end{equation*}
whereas in the non-uniform case,
\begin{equation*}
    x_i - x_{i-1} = h_i, \quad \text{with } h_i \neq h_j \text{ for some } i \neq j.
\end{equation*}
As an illustrative example, one can construct a uniform overlapping decomposition of the domain is obtained by extending each subinterval by a fixed extension distance $\delta > 0$, that is,
\begin{equation*}
    \Omega_1 = (a, x_1 + \delta), \qquad \Omega_i = (x_{i-1} - \delta, x_i + \delta), \qquad \Omega_N = (x_{N-1} - \delta, b),
\end{equation*}
for $i = 2, \ldots, N-1$.

Now consider a finite element subdivision of $\Omega$, consisting of element domains $K_j$, $j = 1, \ldots, N_h$, which are closed subintervals.\footnote{In Ciarlet’s definition of a finite element, the element domain $K$ is a bounded closed set with nonempty interior and piecewise smooth boundary.} These elements form a partition of the closure of $\Omega$, i.e.,
\begin{equation*}
    \overline{\Omega} = \bigcup_{j=1}^{N_{h}}K_{j},
\end{equation*}
where each element is of the form
\begin{equation*}
    K_{j} = [y_{j-1}, y_{j}], \qquad |K_{j}| = h.\footnote{Here, $h$ denotes the size of the finite element $K_j$ and is not the same as the $h$ used above for the spacing of the non-overlapping subdomains.}
\end{equation*}
Assume that each subinterval $\widetilde{\Omega}_i$ is formed by grouping element domains, that is,
\begin{equation*}
\widetilde{\Omega}_i = \operatorname{int}\left(\bigcup_{j\in I_{i}} K_{j} \right),
\end{equation*}
where $I_{i}\subset \{1, \ldots, N_{h}\}$ are pairwise disjoint index sets satisfying
\begin{equation*}
\bigcup_{i=1}^{N} I_i = \{1, \ldots, N_h\}.
\end{equation*}
For each $i$, the overlapping index set is defined by
\begin{equation*}
    I_{i}^{\delta} = \{j \in \{1, \ldots, N_h\} : \operatorname{dist}(K_{j}, \widetilde{\Omega}_{i}) < \delta\}.
\end{equation*}
Then the overlapping subdomain $\Omega_i$ is given by
\begin{equation*}
    \Omega_i = \operatorname{int} \left(\bigcup_{j \in I_i^\delta} K_j \right).
\end{equation*}
We refer to this type of decomposition as a \textbf{mesh-conforming overlapping decomposition} of the domain $\Omega$. 
After presenting these illustrative cases, we proceed to give the general definition of an overlapping decomposition valid for arbitrary finite element partitions. 

In this construction, the extension distance is taken as $\delta = c\, h$, where $h$ denotes the mesh size of the finite element partition and $c \in \mathbb{Z}_{>0}$ specifies the number of neighboring elements to include in the overlap. 
For instance, if we wish to extend the subdomain by only one neighboring element, we take $c = 1$.

\textit{Illustrative example:} 
Consider the bounded domain $\Omega = (0,1)$, and a subdivision into $N_h = 10$ finite element domains, each of size $h = 0.1$:
\begin{equation*}
    K_1 = [0,0.1], \quad K_2 = [0.1,0.2], \quad\ldots, \quad K_{10} = [0.9,1].
\end{equation*}
Let the non-overlapping subdomains be defined by grouping elements:
\begin{align*}
    \widetilde{\Omega}_1 &= \operatorname{int}(K_1 \cup K_2 \cup K_3) = (0,0.3),\\
    \widetilde{\Omega}_2 &= \operatorname{int}(K_4 \cup K_5 \cup K_6) = (0.3,0.6),\\
    \widetilde{\Omega}_3 &= \operatorname{int}(K_7 \cup K_8 \cup K_9 \cup K_{10}) = (0.6,1).
\end{align*}
If we choose $\delta = h$, that is, we extend each subdomain by one neighboring element on each side, then the overlapping subdomains become
\begin{align*}
    \Omega_1 &= \operatorname{int}(K_1 \cup K_2 \cup K_3 \cup K_4) = (0,0.4),\\
    \Omega_2 &= \operatorname{int}(K_3 \cup K_4 \cup K_5 \cup K_6 \cup K_7) = (0.2,0.7),\\
    \Omega_3 &= \operatorname{int}(K_6 \cup K_7 \cup K_8 \cup K_9 \cup K_{10}) = (0.5,1).
\end{align*}
Here, each $\Omega_i$ is an open interval consisting of all elements of the corresponding non-overlapping subdomain together with one neighboring element on each side (whenever available). The overlap, therefore, occurs only along full elements, ensuring that the decomposition respects the finite element mesh.
\item Let $d = 2$, so that the domain $\Omega \subset \mathbb{R}^{2}$ is polygonal. Consider a finite element subdivision of $\Omega$ into element domains $K_j$, $j = 1, \ldots, N_h$. Let $\mathcal{T}_h$ denote a triangulation of $\Omega$, i.e.,  
\begin{equation*}
    \overline{\Omega} = \bigcup_{j=1}^{N_h} K_j,
\end{equation*}  
where each element $K_j \in \mathcal{T}_h$ has diameter
\begin{equation*}
    h_{K_j} := \operatorname{diam}(K_j), \quad h := \max_{K_j \in \mathcal{T}_h} h_{K_j}.
\end{equation*}  
Here, $h$ is the maximal element diameter of the mesh. We assume that $\mathcal{T}_h$ is \emph{shape-regular}, meaning there exists a constant $c > 0$, independent of $h$, such that for all $K \in \mathcal{T}_h$,  
\begin{equation*}
    \frac{h_{K_{j}}}{\rho_{K_{j}}} \leq c,
\end{equation*}  
where $\rho_{K_{j}}$ is the radius of the largest inscribed circle in $K_{j}$. Shape-regularity ensures that the elements are well-proportioned and avoids degenerate triangles.  

Suppose that each non-overlapping subdomain $\widetilde{\Omega}_i$ is obtained by grouping elements:  
\begin{equation*}
    \widetilde{\Omega}_i = \operatorname{int}\left(\bigcup_{j \in I_i} K_j \right),
\end{equation*}  
where the index sets $I_i \subset \{1, \ldots, N_h\}$ are pairwise disjoint and satisfy  
\begin{equation*}
    \bigcup_{i=1}^{N} I_i = \{1, \ldots, N_h\}.
\end{equation*}  
For each $i$, define the overlapping index set
\begin{equation*}
    I_i^\delta = \{ j \in \{1, \ldots, N_h\} : \operatorname{dist}(K_j, \widetilde{\Omega}_i) < \delta \},
\end{equation*}  
so that the corresponding overlapping subdomain is  
\begin{equation*}
    \Omega_i = \operatorname{int} \left(\bigcup_{j \in I_i^\delta} K_j \right).
\end{equation*}  
If the extension is intended to include only a single layer of neighboring elements (vertex neighbors), the extension distance can be chosen as  
\begin{equation*}
    \delta = d^{*} := \min_{K_{j}\cap K_{k} = \emptyset}\operatorname{dist}(K_{j}, K_{k}),
\end{equation*} 
where $d^{*}$ is the minimal positive separation distance between non-neighbouring elements.\footnote{Since $\mathcal{T}_h$ is a conforming triangulation, the distance $d^{*}$ is well-defined.}

\noindent Then the overlapping index set becomes  
\begin{equation*}
    I_i^{d^{*}} = \{ j \in \{1, \ldots, N_h\} : \operatorname{dist}(K_j, \widetilde{\Omega}_i) < d^{*}\},
\end{equation*}  
and the overlapping subdomain is  
\begin{equation*}
    \Omega_i = \operatorname{int} \left(\bigcup_{j \in I_i^{d^{*}}} K_j\right).
\end{equation*}  
This construction can be generalized to an arbitrary number $l$ of extension layers. Let $\Omega_i^0 := \widetilde{\Omega}_i$, and define recursively  
\begin{equation*}
    I_i^{(ld^{*})} = \{ j \in \{1, \ldots, N_h\} : \operatorname{dist}(K_j, \Omega_i^{l-1}) < d^{*}\}, \qquad \Omega_i^l = \operatorname{int} \left(\bigcup_{j \in I_i^{(ld^{*})}} K_j \right), \quad l \geq 1.
\end{equation*}  
This defines the overlapping subdomains layer by layer, with each layer including all elements within distance $d^{*}$ of the previous layer.
% \begin{figure}[H]
%     \centering
%     \includesvg[width=0.25\textwidth]{triext.svg}
%     \caption{One-layer extension of a triangular subdomain $\widetilde{\Omega}_{i}$ on an unstructured triangular mesh.}
%     \label{fig3}
% \end{figure}
As illustrated in Figure~\ref{fig3}, the overlapping subdomain $\Omega_i$ is obtained by including all elements within distance $d^{*}$ of the non-overlapping subdomain $\widetilde{\Omega}_i$, corresponding to a one-layer vertex neighborhood.
\end{enumerate}
We now state the general definition of a mesh-conforming overlapping decomposition.
\begin{definition}{Mesh-Conforming Overlapping Domain Decomposition}{meshdecomp}
    Let $\{K_j\}_{j=1}^{N_{h}}$ be a finite element subdivision of a polytopal domain $\Omega\subset \mathbb{R}^{d}$, and let
    \begin{equation*}
        \widetilde{\Omega}_{i} = \operatorname{int}\left(\bigcup_{j \in I_i} K_j \right),
    \end{equation*}
    be a non-overlapping subdomain of $\Omega$, where the index sets $I_i \subset \{1, \ldots, N_h\}$ are pairwise disjoint. A \textbf{mesh-conforming overlapping decomposition} of $\Omega$ is a finite collection of subdomains $\{\Omega_i\}$ such that
    \begin{equation*}
        \Omega_i = \operatorname{int} \left(\bigcup_{j \in I_i^\delta} K_j \right),
    \end{equation*}
    where the overlapping index set is defined by  
    \begin{equation*}
        I_{i}^{\delta} = \{j \in \{1, \ldots, N_h\} : \operatorname{dist}(K_{j}, \widetilde{\Omega}_{i}) < \delta_{i}\}.
    \end{equation*}
    and $\delta_{i} > 0$ is the extension distance.
\end{definition}
\noindent The definition is equivalent to saying that each overlapping subdomain $\Omega_i$ can be expressed as
\begin{equation*}
    \Omega_{i} = \{x\in\Omega : \operatorname{dist}(x, \widetilde{\Omega}_{i}) < \delta_{i}(x)\},
\end{equation*}
with the \textit{minimal extension condition}
\begin{equation*}
    \operatorname{dist}(\widetilde{\Gamma}_i^{\text{int}}, \Gamma_i^{\text{int}}) \geq \min_{1\leq i \leq N}\delta_{i} = : \delta,
\end{equation*}
where $\delta$ is the minimal extension distance. This equivalence also shows that a mesh-conforming overlapping decomposition is generally an 
\emph{overlapping decomposition with $\delta$-extension}. However, in certain mesh constructions, for example in structured meshes, this can 
correspond to a \emph{uniform overlapping decomposition with $\delta$-extension}, where each subdomain is extended by the same fixed distance.

\vspace{10pt}
\noindent In the paper by Maksymilian Dryja and Olof Widlund~\cite{DryjaWidlund1987}, where an additive Schwarz method for elliptic problems is presented, they choose the local extension for each subdomain, as defined in \ref{meshdecomp}, to be
\begin{equation*}
    \delta_i = \gamma \cdot \operatorname{diam}(\widetilde{\Omega}_{i}) := \gamma \cdot H_{i},
\end{equation*}
where $\gamma \in (0,1)$. This implies that the minimal extension distance is 
\begin{equation*}
    \delta = \min_{1\leq i \leq N}\delta_{i} = \gamma \cdot \min_{1\leq i \leq N}\operatorname{diam}(\widetilde{\Omega}_{i}).
\end{equation*}
Note that if we choose
\begin{equation*}
    \gamma = \dfrac{d^{*}}{\operatorname{diam}(\widetilde{\Omega}_{i})} < 1,
\end{equation*}
then the corresponding overlapping subdomain $\Omega_i$ is obtained by a one-layer extension of the triangular subdomain $\widetilde{\Omega}_i$.